\documentclass[11pt,a4paper]{article}
% preamble.tex — estilo común de todos los apuntes Froth.
% Cada apunte hace % preamble.tex — estilo común de todos los apuntes Froth.
% Cada apunte hace % preamble.tex — estilo común de todos los apuntes Froth.
% Cada apunte hace \input{preamble} justo después de \documentclass.
\usepackage[margin=2.4cm]{geometry}
\usepackage{xcolor}
\definecolor{litblue}{RGB}{31,79,121}
\definecolor{litgreen}{RGB}{27,94,32}
\definecolor{litorange}{RGB}{230,81,0}
\usepackage{parskip}
\usepackage{titlesec}
\titleformat{\section}{\Large\bfseries\color{litblue}}{\thesection}{0.6em}{}
\titleformat{\subsection}{\large\bfseries\color{litblue}}{\thesubsection}{0.6em}{}
\usepackage{tcolorbox}
\usepackage{fancyhdr}
\pagestyle{fancy}
\fancyhf{}
\fancyhead[L]{\small\color{litblue}Froth — Apuntes de ML}
\fancyhead[R]{\small\thepage}
\renewcommand{\headrulewidth}{0.4pt}
\usepackage[colorlinks=true,linkcolor=litblue,urlcolor=litblue]{hyperref}

% Cabecera de cada apunte: \cabecera{numero}{titulo}
\newcommand{\cabecera}[2]{%
  \begin{tcolorbox}[colback=litblue,coltext=white,colframe=litblue,arc=2mm]
    {\Large\bfseries Apunte #1 — #2}\\[3pt]
    {\small Proyecto Froth · aprender Machine Learning construyendo}
  \end{tcolorbox}\vspace{0.8em}}

% Cajas temáticas
\newtcolorbox{concepto}[1]{colback=blue!4,colframe=litblue,title={Concepto: #1},arc=1.5mm}
\newtcolorbox{practica}{colback=green!5,colframe=litgreen,title={Pruébalo tú (teclado en mano)},arc=1.5mm}
\newtcolorbox{ojo}{colback=orange!8,colframe=litorange,title={Ojo / error típico},arc=1.5mm}
\newtcolorbox{resumen}{colback=gray!8,colframe=gray!60,title={En una frase},arc=1.5mm}

\newcommand{\codigo}[1]{\texttt{#1}}
 justo después de \documentclass.
\usepackage[margin=2.4cm]{geometry}
\usepackage{xcolor}
\definecolor{litblue}{RGB}{31,79,121}
\definecolor{litgreen}{RGB}{27,94,32}
\definecolor{litorange}{RGB}{230,81,0}
\usepackage{parskip}
\usepackage{titlesec}
\titleformat{\section}{\Large\bfseries\color{litblue}}{\thesection}{0.6em}{}
\titleformat{\subsection}{\large\bfseries\color{litblue}}{\thesubsection}{0.6em}{}
\usepackage{tcolorbox}
\usepackage{fancyhdr}
\pagestyle{fancy}
\fancyhf{}
\fancyhead[L]{\small\color{litblue}Froth — Apuntes de ML}
\fancyhead[R]{\small\thepage}
\renewcommand{\headrulewidth}{0.4pt}
\usepackage[colorlinks=true,linkcolor=litblue,urlcolor=litblue]{hyperref}

% Cabecera de cada apunte: \cabecera{numero}{titulo}
\newcommand{\cabecera}[2]{%
  \begin{tcolorbox}[colback=litblue,coltext=white,colframe=litblue,arc=2mm]
    {\Large\bfseries Apunte #1 — #2}\\[3pt]
    {\small Proyecto Froth · aprender Machine Learning construyendo}
  \end{tcolorbox}\vspace{0.8em}}

% Cajas temáticas
\newtcolorbox{concepto}[1]{colback=blue!4,colframe=litblue,title={Concepto: #1},arc=1.5mm}
\newtcolorbox{practica}{colback=green!5,colframe=litgreen,title={Pruébalo tú (teclado en mano)},arc=1.5mm}
\newtcolorbox{ojo}{colback=orange!8,colframe=litorange,title={Ojo / error típico},arc=1.5mm}
\newtcolorbox{resumen}{colback=gray!8,colframe=gray!60,title={En una frase},arc=1.5mm}

\newcommand{\codigo}[1]{\texttt{#1}}
 justo después de \documentclass.
\usepackage[margin=2.4cm]{geometry}
\usepackage{xcolor}
\definecolor{litblue}{RGB}{31,79,121}
\definecolor{litgreen}{RGB}{27,94,32}
\definecolor{litorange}{RGB}{230,81,0}
\usepackage{parskip}
\usepackage{titlesec}
\titleformat{\section}{\Large\bfseries\color{litblue}}{\thesection}{0.6em}{}
\titleformat{\subsection}{\large\bfseries\color{litblue}}{\thesubsection}{0.6em}{}
\usepackage{tcolorbox}
\usepackage{fancyhdr}
\pagestyle{fancy}
\fancyhf{}
\fancyhead[L]{\small\color{litblue}Froth — Apuntes de ML}
\fancyhead[R]{\small\thepage}
\renewcommand{\headrulewidth}{0.4pt}
\usepackage[colorlinks=true,linkcolor=litblue,urlcolor=litblue]{hyperref}

% Cabecera de cada apunte: \cabecera{numero}{titulo}
\newcommand{\cabecera}[2]{%
  \begin{tcolorbox}[colback=litblue,coltext=white,colframe=litblue,arc=2mm]
    {\Large\bfseries Apunte #1 — #2}\\[3pt]
    {\small Proyecto Froth · aprender Machine Learning construyendo}
  \end{tcolorbox}\vspace{0.8em}}

% Cajas temáticas
\newtcolorbox{concepto}[1]{colback=blue!4,colframe=litblue,title={Concepto: #1},arc=1.5mm}
\newtcolorbox{practica}{colback=green!5,colframe=litgreen,title={Pruébalo tú (teclado en mano)},arc=1.5mm}
\newtcolorbox{ojo}{colback=orange!8,colframe=litorange,title={Ojo / error típico},arc=1.5mm}
\newtcolorbox{resumen}{colback=gray!8,colframe=gray!60,title={En una frase},arc=1.5mm}

\newcommand{\codigo}[1]{\texttt{#1}}

\begin{document}
\cabecera{35}{La frase centroide: el filtro de relevancia al revés}

\section{La idea (ya la conocías sin saberlo)}

En la Fase 3, el filtro de relevancia medía el coseno de cada paper contra el vector del
TOPIC y \emph{descartaba} lo lejano. La frase centroide usa la MISMA matemática con la
intención opuesta: medir el coseno de cada frase contra el centro del cluster y
\emph{coronar} lo más cercano.

\begin{concepto}{centroide y frase representativa}
El \textbf{centroide} de un cluster es el promedio de los vectores (768D) de sus papers:
su ``centro de masa semántico''. Las frases de los abstracts se vectorizan con el MISMO
modelo (SPECTER), así que viven en el mismo espacio y son comparables. La frase con mayor
coseno al centroide es la que mejor resume de qué va el subtema --- elegida por geometría,
no por opinión.
\end{concepto}

\section{Los resultados reales (corpus Beauvoir, 237 papers)}

\texttt{python -m froth.review} sobre tus datos:

\begin{itemize}
  \item Cluster geología (``ma, melt, compositions''): score \textbf{0.947} para
        \emph{``LCT-type granitic pegmatites represent a major global source of
        lithium...''} --- una definición de libro del subtema.
  \item Cluster USGS (``usgs, old scrap, year''): \emph{``Minerals Yearbook --- These
        annual publications review the mineral industries of the United States...''}
        Las keywords jamás te dijeron que ese cluster son INFORMES ANUALES; la frase sí.
  \item Cluster flotación de lepidolita: \emph{``Efficient flotation of lepidolite is
        crucial for the extraction and recycling of lithium resources...''}
\end{itemize}

Esto resuelve tu pedido de títulos con sentido (7.1b): keywords para clasificar,
\textbf{frase centroide como subtítulo} para entender.

\section{Dos guardas contra trampas}

\begin{itemize}
  \item \textbf{Tope por paper} (\texttt{per\_paper\_cap=2}): sin él, un abstract muy
        ``central'' podría monopolizar el resumen entero del cluster.
  \item El centroide se calcula sobre los vectores de PAPER (los del mapa), no sobre las
        frases: así el resumen es fiel al mismo espacio que el usuario está viendo.
\end{itemize}

\begin{ojo}
Las frases top se parecen al centroide... y por eso se parecen ENTRE SÍ: el ranking por
centroide tiende a la redundancia (tres frases diciendo lo mismo). Ese es exactamente el
problema que ataca MMR en la sub-fase 7.3.
\end{ojo}

\begin{resumen}
Frase centroide = coseno(frase, centro del cluster), misma matemática que el filtro de
relevancia pero para coronar en vez de descartar. En Beauvoir produce definiciones de
libro con scores 0.92-0.95. Su debilidad (redundancia) la corrige MMR (7.3).
\end{resumen}

\end{document}
